% TOP_PROBLEM: {"id": 38, "title": "The Köthe Conjecture", "category_id": 4, "category": {"id": 4, "name": "algebra", "display_name": "Algebra", "description": "Group theory, ring theory, field theory, and algebraic structures.", "slug": "algebra", "order_index": 4, "created_at": "2026-07-31T15:26:25.671Z"}, "status": "open"}
% FIRST_POSTED: null
% ENTRY_KIND: statement_only
% Imported from: batch01.tex; UnsolvedMath ID 38
\documentclass[11pt]{article}
\usepackage[margin=25mm]{geometry}
\usepackage{fontspec}
\setmainfont{Latin Modern Roman}
\usepackage{amsmath,amssymb,mathtools}
\usepackage{unicode-math}
\setmathfont{Latin Modern Math}
\usepackage{xurl,hyperref}
\hypersetup{colorlinks=true,urlcolor=blue}
\setcounter{secnumdepth}{0}
\setlength{\parindent}{0pt}
\setlength{\parskip}{5pt}
\setlength{\emergencystretch}{3em}
\newcommand{\sref}[2]{\hyperlink{source:#1:#2}{[S#2]}}
\newcommand{\cataloguescope}[1]{\paragraph{Recorded target.}#1}
\begin{document}
% UnsolvedMath ID 38; source catalogue code ALG-003
\setcounter{section}{37}
\section{The Köthe Conjecture}\label{38}
{\small\sffamily UnsolvedMath ID 38 · Algebra}
\subsection{Definitions and mathematical statement}

\paragraph{Shared notation.}$\mathbb N=\{1,2,\ldots\}$, and $\mathbb Z,\mathbb Q,\mathbb R,\mathbb C$ denote the integers, rationals, reals and complex numbers. Any other conventions are specified locally.


Let $R$ be any associative ring; neither commutativity nor a multiplicative identity is assumed. A left ideal is an additive subgroup $J\le(R,+)$ with $RJ\subseteq J$; a right ideal has $JR\subseteq J$. A two-sided ideal satisfies both conditions. An element $x$ is nilpotent if $x^m=0$ for some integer $m\ge1$. An ideal is nil if every one of its elements is nilpotent; the exponent may depend on the element, so a nil ideal need not be a nilpotent ideal. \sref{38}{2}
\paragraph{Statement recorded in the supplied source.}
For every associative ring $R$, is the following biconditional true?
\[
 \begin{aligned}
 &R\text{ has no nonzero nil two-sided ideal}\\
 &\qquad\Longleftrightarrow
 R\text{ has no nonzero nil left ideal and no nonzero nil right ideal}.
 \end{aligned}
\]
The nontrivial implication asserts that the absence of nil two-sided ideals forces the absence of nil ideals on either side; the reverse implication is immediate. \sref{38}{2}
\subsection{Short English statement}
Must a ring with no nonzero nil two-sided ideal also have no nonzero nil one-sided ideal? The question concerns arbitrary associative rings, including rings without an identity.
\subsection{Sources}
\begin{description}
\item[\hypertarget{source:38:1}{S1}] ULAM.AI and the UnsolvedMath Contributors, UnsolvedMath: A Curated Collection of Open Mathematics Problems, record 38 (ALG-003). CC BY 4.0. \url{https://huggingface.co/datasets/ulamai/UnsolvedMath}.
\item[\hypertarget{source:38:2}{S2}] Peter Kálnai and Jan Žemlička, Self-Injective von Neumann Regular Rings and Köthe’s Conjecture, arXiv:1912.12159v2 (2020), opening paragraphs and definitions on page 1. \url{https://arxiv.org/pdf/1912.12159}.
\end{description}
\label{end:38}
\end{document}
